\section*{Supplementary results: additional model development}

Nine controlled runs compared matched attention, full fixed scaling and full learned gating using seeds 17, 29 and 43 (Supplementary Table S15; Fig. S4). Internal AUROC means $\pm$ sample standard deviations were $0.9348\pm0.0044$, $0.9398\pm0.0153$ and $0.9353\pm0.0070$, respectively. Corresponding external values were $0.7270\pm0.1035$, $0.7445\pm0.1254$ and $0.7633\pm0.1581$. The gate comparison showed no consistent advantage across seeds, and the configurations differed in parameter count. Historical compound ablations remain separately documented in Supplementary Table S9 and Fig. S3.

A subsequent bounded study completed 34 neural training runs: ten configurations at three initial seeds, followed by two additional seeds for the internally selected candidate and the Hybrid control (Supplementary Tables S16 and S17). The screen included normalized embeddings, attention masking, direct/subspace attention and two-bag training. It reproduced the original matched-attention and Hybrid controls; these repeated controls are not additional independent clinical evidence. Normalized attention ranked first among eight candidates on mean internal AUROC. Its improvement over Hybrid was 0.0011, below the protocol's 0.01 target. Candidate selection was locked before new external predictions, although this external collection had already been examined in earlier development.

Across all five seeds (17, 29, 43, 59 and 71), internal AUROC was $0.9348\pm0.0134$ for Hybrid and $0.9341\pm0.0084$ for normalized attention. External AUROC was $0.8056\pm0.1299$ and $0.8061\pm0.0916$, respectively. Mean external sensitivities at threshold 0.5 were 49.21\% and 55.66\%, with specificities of 82.50\% and 81.25\%. Thus, the historical primary checkpoint's operating point did not characterize the retrained models. Additional-seed results are reported separately in Supplementary Table S17; no favorable seed replaced the primary checkpoint.

Train-only logistic regression on mean UNI2-h features achieved internal/external AUROCs of 0.9314/0.8335 on the reference bags. Averaging normalized patch features before logistic regression gave 0.9247/0.8462. External sensitivities were 45.39\% and 46.05\%, respectively, at specificity 93.75\% for both. These simple controls retained discrimination but showed a different threshold tradeoff. Complete run metrics, full-feature inference and prespecified aggregation analyses accompany Supplementary Tables S16, S17 and S19. The experiments did not establish a stable advantage of the more complex architecture.

